\section{CUDA GPU 并行编程：向量归约代码解析与性能剖析}

\subsection{引言：CPU基准验证}
在编写和测试 GPU 核函数（Kernel）时，验证其输出的正确性是第一步。由于 CPU 串行代码逻辑简单且结果确定，通常先编写一个 CPU 版本的向量归约函数，将其输出作为“Ground Truth（基准真值）”，以此来校验 GPU 核函数的计算结果。

\begin{lstlisting}[caption={CPU端向量归约基准函数}]
float cpu_reduce(const float* input, int n) {
    float sum = 0.0f;
    for (int i = 0; i < n; i++) {
        sum += input[i];
    }
    return sum;
}
\end{lstlisting}

\subsection{CUDA 主函数（Host端）完整实现}
一个标准的 CUDA 应用程序通常包含以下七个核心步骤：分配内存、初始化数据、主机到设备的数据拷贝（H2D）、配置并启动核函数、设备同步、设备到主机的数据拷贝（D2H）、释放内存。


\subsection{内存分配与初始化}
假设向量大小 $N = 1024 \times 1024$（约 100 万元素）。
\begin{itemize}
    \item \textbf{主机端（CPU）}：使用 \lstinline{new} 或 \lstinline{malloc} 分配内存。
    \item \textbf{设备端（GPU）}：使用 \lstinline{cudaMalloc} 分配全局内存。注意计算大小时需乘以 \lstinline{sizeof(float)}。
\end{itemize}

\begin{lstlisting}[caption={内存分配与数据初始化}]
int n = 1024 * 1024;
size_t size = n * sizeof(float);

// 1. 分配内存
float *h_input = new float[n];
float *d_input;
cudaMalloc((void**)&d_input, size);

// 2. 初始化主机数据
for (int i = 0; i < n; i++) {
    h_input[i] = static_cast<float>(i); 
}

// 3. 主机到设备的数据拷贝 (H2D)
cudaMemcpy(d_input, h_input, size, cudaMemcpyHostToDevice);
\end{lstlisting}

\subsection{网格与线程块配置（Grid/Block Sizing）}
为了保证所有元素都被处理，且网格大小（Grid Size）为整数，使用标准的向上取整公式：
\begin{equation}
    \text{gridSize} = \frac{N + \text{blockSize} -1}
{\text{blockSize}}
\end{equation}
假设 \lstinline{blockSize = 256}，则 \lstinline
{gridSize =(1048576 + 255) / 256 = 4096}。

\subsection{多级核函数启动（层级归约）}
由于 CUDA 不支持跨 Block 的全局同步，当向量极大时，单次核函数无法将所有元素归约为 1 个值。需要采用 \textbf{多级核函数启动} 策略，将上一次核函数输出的“部分和（Partial Sums）”作为下一次核函数的输入。


\begin{enumerate}
    \item \textbf{第一级核函数}：处理 $1,048,576$ 个元素，启动 $4096$ 个 Block。输出 $4096$ 个部分和。
    \item \textbf{第二级核函数}：处理 $4096$ 个部分和，启动 $16$ 个 Block（$4096/256$）。输出 $16$ 个部分和。
    \item \textbf{第三级核函数}：处理 $16$ 个部分和，启动 $1$ 个 Block。输出 $1$ 个最终总和。
\end{enumerate}

\begin{lstlisting}[caption={多级核函数启动与同步}]
int blockSize = 256;

// --- 第一级归约 ---
int gridSize1 = (n + blockSize - 1) / blockSize; // 4096
reduce_in_place<<<gridSize1, blockSize>>>(d_input, n);
cudaDeviceSynchronize(); // 必须同步，确保部分和已写入全局内存

// --- 第二级归约 ---
int n2 = gridSize1; // 4096
int gridSize2 = (n2 + blockSize - 1) / blockSize; // 16
reduce_in_place<<<gridSize2, blockSize>>>(d_input, n2);
cudaDeviceSynchronize();

// --- 第三级归约 ---
int n3 = gridSize2; // 16
int gridSize3 = 1; // 只需1个Block
reduce_in_place<<<gridSize3, blockSize>>>(d_input, n3);
cudaDeviceSynchronize();
\end{lstlisting}

\subsection{结果回传、验证与资源释放}
最终结果存储在 \lstinline{d_input[0]} 中。只需将这 4 个字节（1 个 float）拷回主机。

\begin{lstlisting}[caption={结果回传与内存释放}]
// 仅拷贝第一个元素（最终结果）回主机
cudaMemcpy(h_input, d_input, sizeof(float), cudaMemcpyDeviceToHost);

std::cout << "GPU Output: " << h_input[0] << std::endl;
std::cout << "CPU Output: " << cpu_reduce(h_input, n) << std::endl;

// 释放内存
cudaFree(d_input);
delete[] h_input;
\end{lstlisting}

\begin{tcolorbox}[colback=yellow!5!white,colframe=yellow!75!black,title=浮点精度差异说明]
运行结果可能会显示 GPU 输出（如 5.49...e11）与 CPU 输出在最低有效位上存在微小差异。这是因为浮点数加法
不满足结合律，且 GPU 和 CPU 的累加顺序、硬件舍入模式不同。若需完全一致，可将数据类型改为 \lstinline
{double}，或使用 C++ 模板（Templates）来统一管理数据类型。
\end{tcolorbox}

\subsection{Nsight Compute 性能剖析 (Profiling)}
代码跑通后，使用 NVIDIA Nsight Compute（NCU）对应用程序进行性能分析，以寻找瓶颈。

\subsection{核函数执行时间分析（Duration）}
在 NCU 的 Summary 页面，可以看到三个同名核函数的执行时间：
\begin{itemize}
    \item \textbf{核函数 1（4096 Blocks）}：约 284 $\mu s$。耗时最长，因为处理了 100 万个元素。
    \item \textbf{核函数 2（16 Blocks）}：约 15 $\mu s$。
    \item \textbf{核函数 3（1 Block）}：约 13 $\mu s$。
\end{itemize}

\subsection{SM 占用与波次 (Waves) 概念}
GPU 的计算单元由多个 SM 组成。假设当前 GPU 有 38 个 SM：
\begin{itemize}
    \item \textbf{核函数 1}：4096 个 Block 分配给 38 个 SM。每个 SM 需要执行约 $4096 / 38 \approx 108$ 个 Block。这被称为 \textbf{多波次（Multi-wave）} 执行，SM 处于高负载状态。
    \item \textbf{核函数 2}：仅有 16 个 Block。这意味着只有 16 个 SM 各执行 1 个 Block，剩余 22 个 SM 完全闲置。这导致了严重的硬件利用率下降（Tail Effect / 尾部效应）。
\end{itemize}

\subsection{计算与内存吞吐量}
\begin{itemize}
    \item \textbf{SM 利用率 (Compute Throughput)}：Kernel 1 的计算利用率约为 74\%。
    \item \textbf{内存吞吐量 (Memory Throughput)}：仅为 9\% 左右。
    \item \textbf{结论}：当前版本是 \textbf{计算受限（Compute-bound）}或受限于指令开销，而非内存带宽受限。大量的线程同步（\lstinline{__syncthreads}）和条件分支（\lstinline{if} 过滤器）消耗了大量时钟周期。
\end{itemize}

\subsection{缓存命中率 (Cache Hit Rate) 的对比}
在之前的“向量加法”中，每个元素只被读取一次（Streaming），L1/L2 缓存命中率几乎为 0\%。但在“向量归约”中，数据存在强烈的 \textbf{时间局部性（Temporal Locality）}：
\begin{itemize}
    \item 第 0 步：读取 index 0 和 1。
    \item 第 1 步：再次读取 index 0（此时包含 0+1 的结果）和 2。
    \item \textbf{结果}：同一个内存地址被反复读取和更新。因此，NCU 显示 \textbf{L1 缓存命中率高达 90\%，L2 缓存命中率约 80\%}。这极大地缓解了全局内存带宽的压力。
\end{itemize}

\subsection{总结与后续优化方向}
本章完成了向量归约的完整 Host/Device 代码编写，并学会了使用 Nsight Compute 分析性能指标。当前的“仅使用全局内存 + 跨步访问”版本存在以下明显缺陷，将在后续章节中优化：
\begin{enumerate}
    \item \textbf{非合并访存（Non-coalesced Access）}：随着 Stride 增大，相邻线程访问的内存地址跨度越来越大，导致内存事务（Memory Transactions）浪费。
    \item \textbf{线程发散（Thread Divergence）}：\lstinline{if (tid \% (2*stride) == 0)} 导致 Warp 内一半线程闲置，产生严重的分支发散。
    \item \textbf{尾部效应（Tail Effect）}：最后几级 Kernel 启动的 Block 数极少，导致大量 SM 闲置。
\end{enumerate}
\textbf{下一步计划}：引入 \textbf{共享内存（Shared Memory）} 来消除全局内存的重复访问，
并采用 \textbf{交错寻址（Interleaved Addressing）} 或 \textbf{Warp Shuffle} 指令来彻底解决线程发散和同步开销问题。
